% GENERATED FILE. Edit canonical lessons, then run npm run generate:books.
% canonical-source-hash: fnv1a64:e2eb04fce57a0684
% canonical-lessons: GE-C477-tapfer, GE-C477-feige, GE-C477-feigling, GE-C477-leichtsinnig, GE-C477-trauen

\chapter{Tapfer oder feige}
\label{ch:ge-tapfer-oder-feige}
\hlchaptermodality{477}

\begin{chapteropening}
\textbf{By the end of this chapter:} \emph{I can talk about courage and cowardice.}

Complete the last lesson of chapter 477: I can talk about courage and cowardice.
\end{chapteropening}

\section[tapfer]{tapfer — brave}

\label{lesson:GE-C477-tapfer}



\begin{quote}
Before the new one: say the German for irritable, then the German for sensitive.
\end{quote}

\subsection*{You'll want to know: tapfer}
\textbf{tapfer} — \textquotedblleft{}brave\textquotedblright{}.

\begin{quote}
Das Kind war beim Arzt sehr tapfer.
\end{quote}

\begin{quote}
\emph{The child was very brave at the doctor's.}
\end{quote}

The noun: \emph{die Tapferkeit} — bravery.

\begin{cousinweb}
\emph{tapfer} is inherited from Middle High German \emph{tapfer} and Old High German \emph{tapfar}, from Proto-Germanic \emph{dapraz} \textquotedblleft{}stout, solid, heavy\textquotedblright{} (EtymDB); the sense moved from firm to brave.
\end{cousinweb}

\begin{grammarlens}[title={What you've built}]
One more word for telling, explaining and describing at B1.
\end{grammarlens}

\subsection*{Your turn}
\begin{itemize}
\raggedright
  \item \emph{Say it:} \emph{tapfer}
  \item \emph{Say it:} \emph{tapfer}, once more
  \item \emph{Say it:} a sentence of your own with \emph{tapfer}
  \item \emph{Recall:} say the German for irritable, then the German for sensitive, then say \emph{tapfer} again
\end{itemize}

\begin{tcolorbox}[breakable,colback=teal!4,colframe=teal!35!black,title={Before you move on}]
What does tapfer mean? (\textquotedblleft{}Brave\textquotedblright{}.) Say it once more.
\end{tcolorbox}
\section[feige]{feige — cowardly}

\label{lesson:GE-C477-feige}



\begin{quote}
Before the new one: say the German for sensitive, then the German for brave.
\end{quote}

\subsection*{You'll want to know: feige}
\textbf{feige} — \textquotedblleft{}cowardly\textquotedblright{}.

\begin{quote}
Es war feige, einfach wegzulaufen.
\end{quote}

\begin{quote}
\emph{It was cowardly to simply run away.}
\end{quote}

Not to be confused with \emph{die Feige} — the fig.

\begin{cousinweb}
\emph{feige} is inherited from Old High German \emph{feigi} \textquotedblleft{}doomed to die\textquotedblright{}, from Proto-Germanic \emph{faigijaz} (EtymDB); English \emph{fey} is a cousin.
\end{cousinweb}

\begin{grammarlens}[title={What you've built}]
One more word for telling, explaining and describing at B1.
\end{grammarlens}

\subsection*{Your turn}
\begin{itemize}
\raggedright
  \item \emph{Say it:} \emph{feige}
  \item \emph{Say it:} \emph{feige}, once more
  \item \emph{Say it:} a sentence of your own with \emph{feige}
  \item \emph{Recall:} say the German for sensitive, then the German for brave, then say \emph{feige} again
\end{itemize}

\begin{tcolorbox}[breakable,colback=teal!4,colframe=teal!35!black,title={Before you move on}]
What does feige mean? (\textquotedblleft{}Cowardly\textquotedblright{}.) Say it once more.
\end{tcolorbox}
\section[der Feigling]{der Feigling — coward}

\label{lesson:GE-C477-feigling}



\begin{quote}
Before the new one: say the German for brave, then the German for cowardly.
\end{quote}

\subsection*{You'll want to know: der Feigling}
\textbf{der Feigling} — \textquotedblleft{}coward\textquotedblright{}.

\begin{quote}
Im Film ist der Held am Anfang ein Feigling.
\end{quote}

\begin{quote}
\emph{In the film the hero is a coward at first.}
\end{quote}

Plural: \emph{die Feiglinge}.

\begin{cousinweb}
\emph{Feigling} is built on \emph{feige} \textquotedblleft{}cowardly\textquotedblright{} with \emph{-ling}; \emph{feige} is inherited from Old High German \emph{feigi}, from Proto-Germanic \emph{faigijaz} (EtymDB).
\end{cousinweb}

\begin{grammarlens}[title={What you've built}]
One more word for telling, explaining and describing at B1.
\end{grammarlens}

\subsection*{Your turn}
\begin{itemize}
\raggedright
  \item \emph{Say it:} \emph{der Feigling}
  \item \emph{Say it:} \emph{der Feigling}, once more
  \item \emph{Say it:} a sentence of your own with \emph{der Feigling}
  \item \emph{Recall:} say the German for brave, then the German for cowardly, then say \emph{der Feigling} again
\end{itemize}

\begin{tcolorbox}[breakable,colback=teal!4,colframe=teal!35!black,title={Before you move on}]
What does der Feigling mean? (\textquotedblleft{}Coward\textquotedblright{}.) Say it once more.
\end{tcolorbox}
\section[leichtsinnig]{leichtsinnig — reckless, careless}

\label{lesson:GE-C477-leichtsinnig}



\begin{quote}
Before the new one: say the German for cowardly, then the German for coward.
\end{quote}

\subsection*{You'll want to know: leichtsinnig}
\textbf{leichtsinnig} — \textquotedblleft{}reckless, careless\textquotedblright{}.

\begin{quote}
Ohne Helm Fahrrad zu fahren ist leichtsinnig.
\end{quote}

\begin{quote}
\emph{Riding a bike without a helmet is reckless.}
\end{quote}

The noun: \emph{der Leichtsinn} — recklessness.

\begin{cousinweb}
\emph{leichtsinnig} is built on \emph{leicht} \textquotedblleft{}light\textquotedblright{} and \emph{Sinn} \textquotedblleft{}mind, sense\textquotedblright{}, which is inherited from Old High German \emph{sin}, from Proto-Germanic \emph{sinnaz} (EtymDB): a light mind.
\end{cousinweb}

\begin{grammarlens}[title={What you've built}]
One more word for telling, explaining and describing at B1.
\end{grammarlens}

\subsection*{Your turn}
\begin{itemize}
\raggedright
  \item \emph{Say it:} \emph{leichtsinnig}
  \item \emph{Say it:} \emph{leichtsinnig}, once more
  \item \emph{Say it:} a sentence of your own with \emph{leichtsinnig}
  \item \emph{Recall:} say the German for cowardly, then the German for coward, then say \emph{leichtsinnig} again
\end{itemize}

\begin{tcolorbox}[breakable,colback=teal!4,colframe=teal!35!black,title={Before you move on}]
What does leichtsinnig mean? (\textquotedblleft{}Reckless, careless\textquotedblright{}.) Say it once more.
\end{tcolorbox}
\section[trauen]{trauen — to dare (sich trauen); to trust}

\label{lesson:GE-C477-trauen}



\begin{quote}
Before the new one: say the German for coward, then the German for reckless.
\end{quote}

\subsection*{You'll want to know: trauen}
\textbf{trauen} — \textquotedblleft{}to dare (sich trauen); to trust\textquotedblright{}.

\begin{quote}
Er traut sich nicht, vom Sprungturm zu springen.
\end{quote}

\begin{quote}
\emph{He doesn't dare to jump off the diving tower.}
\end{quote}

Reflexive for to dare: \emph{sich trauen}. With a dative it means to trust: \emph{Ich traue ihm nicht.}

\begin{cousinweb}
\emph{trauen} is inherited from Old High German \emph{trūwēn} \textquotedblleft{}to trust, to believe\textquotedblright{} (EtymDB), a Germanic verb related to \emph{treu} \textquotedblleft{}faithful\textquotedblright{}.
\end{cousinweb}

\begin{grammarlens}[title={What you've built}]
One more word for telling, explaining and describing at B1.
\end{grammarlens}

\subsection*{Your turn}
\begin{itemize}
\raggedright
  \item \emph{Say it:} \emph{trauen}
  \item \emph{Say it:} \emph{trauen}, once more
  \item \emph{Say it:} a sentence of your own with \emph{trauen}
  \item \emph{Recall:} say the German for coward, then the German for reckless, then say \emph{trauen} again
\end{itemize}

\begin{tcolorbox}[breakable,colback=teal!4,colframe=teal!35!black,title={Before you move on}]
What does trauen mean? (\textquotedblleft{}To dare (sich trauen); to trust\textquotedblright{}.) Say it once more.
\end{tcolorbox}
